\documentclass[10pt,a4paper]{amsart}
\usepackage[margin=19mm]{geometry}
\usepackage{amsmath,amssymb,mathtools}
\usepackage[scr=rsfs,cal=euler]{mathalfa}
\usepackage{xfrac}
\usepackage[unicode,hidelinks,bookmarks=false]{hyperref}
\newcommand{\ie}{i.e.\ }
\newcommand{\cf}{cf.\ }
\newcommand{\resp}{respectively\ }
\newcommand{\Subsection}{\subsection}
\providecommand{\nobreakdash}{\nobreak\hbox{-}}
\setlength{\emergencystretch}{2em}
\multlinegap0pt
\usepackage{bm}
\usepackage{bbm}
\usepackage{dsfont}
\usepackage{xspace}
\usepackage{cancel}
\usepackage{mathtools}
\usepackage{amsthm}
\makeatletter
\@ifundefined{thm}{\newtheorem{thm}{Thm}}{}
\@ifundefined{lemma}{\newtheorem{lemma}[thm]{Lemma}}{}
\makeatother
\newcommand{\ddc}{dd^c}
\newcommand{\tr}{\operatorname{tr}}
\title[The $J$-equation on K\"ahler manifolds under a smooth bounda]{The $J$-equation on K\"ahler manifolds under a smooth boundary cone condition}
\author{Junbang Liu}
\date{Source: arXiv:2608.23747v1 (CC BY 4.0)}
\begin{document}
\maketitle

\noindent\emph{Source and licence.} The $J$-equation on K\"ahler manifolds under a smooth boundary cone condition. Version arXiv:2608.23747v1 (CC BY 4.0), distributed under the Creative Commons Attribution 4.0 International licence, as recorded in the arXiv licence metadata for this exact version and shown on the abstract page https://arxiv.org/abs/2608.23747v1. Full rights notice: licenses/arxiv-2608.23747-CC-BY-4.0.txt.\par\smallskip

\noindent\textit{Editorial scope.} Counted unit: the Lemma whose source label is \texttt{lem:weighted-trace} in the article (source0.tex lines 1779--1798, source label \texttt{lem:weighted-trace}), delivered together with the article's own complete original proof of it (source0.tex lines 1801--2117, one \texttt{proof} environment of 317 lines). No mathematics is written, repaired or re-derived: statement and proof are the article's own text line for line. Required context retained uncounted: eq:B-weighted-lower (lines 1564--1571); eq:Theta-hatC-w (lines 1723--1725); lem:corrected-reference (lines 1726--1743); eq:B-weighted-hatC (lines 1737--1740). Every definition, display and label of those lines is retained unchanged. Omitted from this package and not redistributed: 1--1563, 1572--1722, 1741--1778, 1799--1800, 2118--3223. Nothing inside a retained excerpt is omitted or altered: every retained body is delivered exactly as the article's lines are written, so no omission receipt of any kind accompanies this package. Citations: the retained excerpts carry no citation command at all, so no bibliography entry of the article is transcribed and no reference is redistributed. Reference adaptation: none was needed and none was made; every reference inside the delivered unit resolves inside it, so no unresolved reference or citation is produced. Printed slips kept literal: no printed word, symbol, hypothesis or number of the counted statement or of its proof was altered. Layout and class: the article's own class is \texttt{amsart} with packages including amsthm, amsmath, amssymb, amscd, graphics, latexsym, stmaryrd, empheq, xcolor, hyperref, enumitem, biblatex; the wrapper is the lane's pinned \texttt{amsart} editorial wrapper at 10pt on A4, which restates the theorem-like environments the retained text needs and adds the article's own macro definitions that the retained text needs (\texttt{\textbackslash ddc}, \texttt{\textbackslash tr}), with extra wrapper packages (bm, bbm, dsfont, xspace, cancel, mathtools). The delivered numbering is the unit's own. Assets: no retained span contains a figure environment, a graphics-inclusion command, a PDF-page inclusion, a table or a TikZ picture; no image file is copied into this package, the package declares no asset dependency and no image file is redistributed with it. Licence: version arXiv:2608.23747v1 of this article is distributed under the Creative Commons Attribution 4.0 International licence, as recorded in the arXiv licence metadata for this exact version and shown on the abstract page https://arxiv.org/abs/2608.23747v1. A full rights notice is delivered at licenses/arxiv-2608.23747-CC-BY-4.0.txt. Characters: every retained excerpt is pure ASCII, so the delivered engine, XeLaTeX with its default Latin Modern text font, reports no missing character.
\begin{align}
 C_0-\varepsilon_0\theta_{E}&\geq c_0\kappa,
 \label{eq:C-exceptional-positive}\\
 B_0-\varepsilon_0\theta_{E}&\geq c_0\kappa,
 \label{eq:B-exceptional-positive}\\
 B_0&\geq c_0|s_E|_{h}^{2\gamma}\kappa.
 \label{eq:B-weighted-lower}
\end{align}

\begin{equation}\label{eq:Theta-hatC-w}
 \Theta_t=\widehat C_t+\ddc w_t.
\end{equation}

\begin{lemma}
\label{lem:corrected-reference}
The number $\varepsilon$ can be chosen independently of $t$ so that
there are constants $c_1,C_1,\eta_1>0$ satisfying
\begin{equation}\label{eq:hatC-uniform}
 c_1\kappa\leq\widehat C_t\leq C_1\kappa,
 \qquad
 P_{B_t}(\widehat C_t)\leq d_t-\eta_1.
\end{equation}
Furthermore, after changing the constant in
\eqref{eq:B-weighted-lower},
\begin{equation}\label{eq:B-weighted-hatC}
 B_t\geq c_1|s_E|_{h}^{2\gamma}
 \widehat C_t.
\end{equation}
All $C^2$-norms of $B_t$ and $\widehat C_t$, measured with
respect to $\kappa$, are bounded independently of $t$.
\end{lemma}

\begin{lemma}
\label{lem:weighted-trace}

There is a constant $C>0$, independent of $A$ and $t$, with the
following property. For every $A>0$, at a maximum point in
$\widetilde Y\setminus E$ of
\[
 Q_t=\log\Lambda_t-Aw_t,
\]
either
\begin{equation}\label{eq:weighted-alternative}
 \Lambda_t\leq
 C|s_E|_{h}^{-4\gamma},
\end{equation}
or
\begin{equation}\label{eq:maximum-inequality}
 0\leq-\mathcal L_tQ_t
 \leq C(1+\mathcal S_t)+A(d_t-\mathcal S_t).
\end{equation}
\end{lemma}

\begin{proof}
Write $g=\widehat C_t$, $\widetilde g=\Theta_t$, and
$\theta=B_t$, set $F^{i\bar j}:=\mathcal F_t^{i\bar j}$ and
$\mathcal L:=\mathcal L_t$, and
suppress $t$ from $d_t,\Lambda_t$, and $\mathcal S_t$. We write $|s|$ for
$|s_E|_{h}$.  Fix the maximum point and take
holomorphic coordinates there such that $g_{i\bar j}=\delta_{ij}$,
the first derivatives of $g$ vanish, and $\widetilde g$ is
diagonal.

We first differentiate the equation
$\widetilde g^{p\bar q}\theta_{p\bar q}=d$.  The inverse-matrix
identity $
 (\widetilde g^{p\bar q})_i
 =-\widetilde g^{p\bar b}(\widetilde g_{a\bar b})_i
   \widetilde g^{a\bar q}$
gives
\begin{equation}\label{eq:first-differentiated-J}
 0=-\widetilde g^{r\bar q}\widetilde g^{p\bar s}
 (\widetilde g_{r\bar s})_i\theta_{p\bar q}
 +\widetilde g^{p\bar q}(\theta_{p\bar q})_i.
\end{equation}
Differentiating once more and contracting the differentiating indices
with $g^{i\bar j}$ gives
\begin{align}
0=g^{i\bar j}\Bigl\{&
 \Bigl[\widetilde g^{r\bar q}\widetilde g^{p\bar s}
       (\widetilde g_{r\bar s})_{i\bar j}
 -\widetilde g^{r\bar q}\widetilde g^{p\bar b}
       \widetilde g^{a\bar s}
       (\widetilde g_{a\bar b})_{\bar j}
       (\widetilde g_{r\bar s})_i \notag\\
&\quad
 -\widetilde g^{r\bar b}\widetilde g^{a\bar q}
       \widetilde g^{p\bar s}
       (\widetilde g_{a\bar b})_{\bar j}
       (\widetilde g_{r\bar s})_i\Bigr]\theta_{p\bar q}
 \notag\\
&+2\operatorname{Re}\!\left(
 \widetilde g^{p\bar\ell}\widetilde g^{k\bar q}
 (\widetilde g_{k\bar\ell})_i
 (\theta_{p\bar q})_{\bar j}\right)
 -\widetilde g^{p\bar q}(\theta_{p\bar q})_{i\bar j}
 \Bigr\}.\label{eq:second-differentiated-J}
\end{align}


At the same point,
\[
 \Lambda_k=g^{p\bar q}(\widetilde g_{p\bar q})_k,
 \qquad
 \Lambda_{k\bar\ell}
 =R^{p\bar q}{}_{k\bar\ell}(g)\widetilde g_{p\bar q}
  +g^{p\bar q}(\widetilde g_{p\bar q})_{k\bar\ell}.
\]
Indeed, $(g^{p\bar q})_k=0$ in $g$-normal coordinates and, with
our curvature convention,
$(g^{p\bar q})_{k\bar\ell}=R^{p\bar q}{}_{k\bar\ell}(g)$.
Insert these formulas and add $\Lambda^{-1}$ times the zero identity
\eqref{eq:second-differentiated-J} to
\[
 -\mathcal L\log\Lambda
 =-\frac{F^{k\bar\ell}\Lambda_{k\bar\ell}}{\Lambda}
  +\frac{F^{k\bar\ell}\Lambda_k\Lambda_{\bar\ell}}{\Lambda^2}.
\]
The term
$\Lambda^{-1}g^{i\bar j}F^{r\bar s}
(\widetilde g_{r\bar s})_{i\bar j}$ cancels
$-\Lambda^{-1}F^{k\bar\ell}g^{p\bar q}
(\widetilde g_{p\bar q})_{k\bar\ell}$.  More explicitly, a local
potential for the K\"ahler form $\widetilde g$ gives
$(\widetilde g_{r\bar s})_{i\bar j}
=(\widetilde g_{i\bar j})_{r\bar s}$, and hence
$g^{i\bar j}F^{r\bar s}(\widetilde g_{r\bar s})_{i\bar j}
=F^{r\bar s}g^{i\bar j}(\widetilde g_{i\bar j})_{r\bar s}$;
renaming $(r,s,i,j)$ as $(k,\ell,p,q)$ gives exactly the other
fourth-order contraction.  What remains is
\begin{align}
-\mathcal L\log\Lambda={}&
-\frac1\Lambda g^{i\bar j}\widetilde g^{r\bar q}
 \widetilde g^{p\bar b}\widetilde g^{a\bar s}
 (\widetilde g_{a\bar b})_{\bar j}
 (\widetilde g_{r\bar s})_i\theta_{p\bar q}\notag\\
&-\frac1\Lambda g^{i\bar j}\widetilde g^{r\bar b}
 \widetilde g^{a\bar q}\widetilde g^{p\bar s}
 (\widetilde g_{a\bar b})_{\bar j}
 (\widetilde g_{r\bar s})_i\theta_{p\bar q}\notag\\
&+\frac2\Lambda\operatorname{Re}\!\left(
 g^{i\bar j}\widetilde g^{p\bar\ell}
 \widetilde g^{k\bar q}(\widetilde g_{k\bar\ell})_i
 (\theta_{p\bar q})_{\bar j}\right)
 -\frac1\Lambda g^{i\bar j}\widetilde g^{p\bar q}
 (\theta_{p\bar q})_{i\bar j}\notag\\
&-\frac1\Lambda F^{k\bar\ell}
 R^{p\bar q}{}_{k\bar\ell}(g)\widetilde g_{p\bar q}
 +\frac1{\Lambda^2}F^{k\bar\ell}
 \Lambda_k\Lambda_{\bar\ell}.
\label{eq:expanded-trace-calculation}
\end{align}
Put
$\widetilde g_{i\bar j}=\lambda_i\delta_{ij}$.  Since
$\widetilde g$ is K\"ahler,
$(\widetilde g_{r\bar a})_i=(\widetilde g_{i\bar a})_r$.
Direct substitution of $g^{i\bar j}=\delta_{ij}$ and
$\widetilde g^{i\bar j}=\lambda_i^{-1}\delta_{ij}$ shows that the
first positive quadratic expression occurring with a minus sign in
\eqref{eq:expanded-trace-calculation} is
\[
 \sum_{i,a,p,r}
 \frac{\theta_{p\bar r}}{\lambda_a\lambda_p\lambda_r}
 (\widetilde g_{r\bar a})_i
 \overline{(\widetilde g_{p\bar a})_i}.
\]
In the same coordinates,
$F^{k\bar\ell}=\theta_{\ell\bar k}/(\lambda_k\lambda_\ell)$ and
\[
 \Lambda_k=\sum_a(\widetilde g_{a\bar a})_k
 =\sum_a(\widetilde g_{k\bar a})_a.
\]
For each $a$, define the vector $\xi_a$ by
$(\xi_a)_k=(\widetilde g_{k\bar a})_a/\lambda_k$. The weighted
Cauchy--Schwarz inequality for the positive Hermitian matrix $\theta$
is
\[
 \left(\sum_a\lambda_a\right)
 \left(\sum_a\frac{|\xi_a|_\theta^2}{\lambda_a}\right)
 \geq\left|\sum_a\xi_a\right|_\theta^2.
\]
It gives
\begin{align}
 F^{k\bar\ell}\Lambda_k\Lambda_{\bar\ell}
 &=\sum_{k,\ell}\theta_{\ell\bar k}
 \left(\sum_a\frac{(\widetilde g_{k\bar a})_a}{\lambda_k}\right)
 \overline{\left(\sum_b
 \frac{(\widetilde g_{\ell\bar b})_b}{\lambda_\ell}\right)}
 \notag\\
 &\leq\left(\sum_a\lambda_a\right)
 \sum_a\frac1{\lambda_a}\sum_{k,\ell}
 \theta_{\ell\bar k}
 \frac{(\widetilde g_{k\bar a})_a}{\lambda_k}
 \overline{\frac{(\widetilde g_{\ell\bar a})_a}{\lambda_\ell}}
 \notag\\
 &\leq\Lambda\sum_{i,a,p,r}
 \frac{\theta_{p\bar r}}{\lambda_a\lambda_p\lambda_r}
 (\widetilde g_{r\bar a})_i
 \overline{(\widetilde g_{p\bar a})_i}.
\label{eq:trace-gradient-CS}
\end{align}
Here the last inequality holds because the
last sum contains all the terms with $i=a$.  Thus the first negative
quadratic term in \eqref{eq:expanded-trace-calculation} absorbs its
final gradient term.

The second positive quadratic expression occurring with a minus sign
in \eqref{eq:expanded-trace-calculation} becomes
\[
 \sum_{i,r,p,a}\frac{\theta_{p\bar a}}
 {\lambda_r\lambda_p\lambda_a}
 (\widetilde g_{r\bar p})_i
 \overline{(\widetilde g_{r\bar a})_i},
\]
whereas the mixed term becomes
\[
 2\operatorname{Re}\sum_{i,r,p}
 \frac{(\widetilde g_{r\bar p})_i}{\lambda_r\lambda_p}
 (\theta_{p\bar r})_{\bar i}.
\]
For every fixed $i,r$, complete the square.  Written out completely,
the resulting nonnegative quantity is
\begin{align*}
0\leq\sum_{i,r}\frac1{\lambda_r}\Biggl\{&
 \sum_{p,a}\theta_{p\bar a}
 \frac{(\widetilde g_{r\bar p})_i}{\lambda_p}
 \overline{\frac{(\widetilde g_{r\bar a})_i}{\lambda_a}}
 -2\operatorname{Re}\sum_p
 \frac{(\widetilde g_{r\bar p})_i}{\lambda_p}
 (\theta_{p\bar r})_{\bar i}\\
&+\sum_{p,a}(\theta_{r\bar p})_i\theta^{p\bar a}
 (\theta_{a\bar r})_{\bar i}\Biggr\}.
\end{align*}
Indeed, the expression in braces is the square, with respect to
$\theta$, of the column whose $p$-th component is
\[
 \frac{\overline{(\widetilde g_{r\bar p})_i}}{\lambda_p}
 -\sum_a\theta^{p\bar a}(\theta_{a\bar r})_{\bar i}.
\]
Rearranging the preceding nonnegative-square identity gives
\begin{align}
&-\sum_{i,r,p,a}\frac{\theta_{p\bar a}}
 {\lambda_r\lambda_p\lambda_a}
 (\widetilde g_{r\bar p})_i
 \overline{(\widetilde g_{r\bar a})_i}
 +2\operatorname{Re}\sum_{i,r,p}
 \frac{(\widetilde g_{r\bar p})_i}{\lambda_r\lambda_p}
 (\theta_{p\bar r})_{\bar i}
 \leq
 \sum_{i,r,p,a}\frac1{\lambda_r}
 (\theta_{r\bar p})_i\theta^{p\bar a}
 (\theta_{a\bar r})_{\bar i}.
\label{eq:mixed-square}
\end{align}
  Combining \eqref{eq:trace-gradient-CS} and
\eqref{eq:mixed-square} with
\eqref{eq:expanded-trace-calculation} therefore yields, at the chosen
normal-coordinate point,
\begin{align}
-\mathcal L\log\Lambda\leq{}&
-\frac1\Lambda g^{i\bar j}\widetilde g^{p\bar q}
 (\theta_{p\bar q})_{i\bar j}
 +\frac1\Lambda\sum_{i,r,p,a}\frac1{\lambda_r}
 (\theta_{r\bar p})_i\theta^{p\bar a}
 (\theta_{a\bar r})_{\bar i}\notag\\
&-\frac1\Lambda F^{k\bar\ell}
 R^{p\bar q}{}_{k\bar\ell}(g)\widetilde g_{p\bar q}.
\label{eq:surviving-coefficients}
\end{align}

We estimate the three terms in
\eqref{eq:surviving-coefficients}.
Lemma~\ref{lem:corrected-reference} gives uniform
$C^2$-bounds for $g$ and $\theta$.  At a $g$-normal point,
first coordinate derivatives are covariant derivatives, while
$(\theta_{p\bar q})_{i\bar j}$ differs from the corresponding second
covariant derivative only by curvature components of $g$ contracted
with $\theta$.  Thus
$|(\theta_{p\bar q})_i|+|(\theta_{p\bar q})_{i\bar j}|\leq C$ in
the above coordinates.  In particular,
\[
 \left|g^{i\bar j}\widetilde g^{p\bar q}
 (\theta_{p\bar q})_{i\bar j}\right|
 =\left|\sum_{i,p}\frac{(\theta_{p\bar p})_{i\bar i}}
 {\lambda_p}\right|
 \leq C\sum_p\frac1{\lambda_p}
 =C\tr_{\widetilde g}g.
\]
For the quadratic coefficient term, positivity of $\theta^{-1}$
gives
\[
 |\theta^{p\bar a}|
 \leq(\theta^{p\bar p}\theta^{a\bar a})^{1/2},
 \qquad
 \sum_{p,a}|\theta^{p\bar a}|
 \leq m\sum_p\theta^{p\bar p}=m\tr_\theta g.
\]
All first derivatives of $\theta$ are uniformly bounded in these
coordinates.  Therefore
\[
 \sum_{i,r,p,a}\frac1{\lambda_r}
 |(\theta_{r\bar p})_i|\,|\theta^{p\bar a}|\,
 |(\theta_{a\bar r})_{\bar i}|
 \leq C\left(\sum_r\frac1{\lambda_r}\right)\tr_\theta g
 =C\tr_{\widetilde g}g\,\tr_\theta g.
\]
Finally, the curvature contraction can also be seen directly in these
coordinates.  For every vector $\xi$,
\[
 \left|R^{p\bar q}{}_{k\bar\ell}(g)
 \widetilde g_{p\bar q}\xi^k\overline{\xi^\ell}\right|
 =\left|\sum_p\lambda_pR^{p\bar p}{}_{k\bar\ell}(g)
 \xi^k\overline{\xi^\ell}\right|
 \leq C\Lambda|\xi|_g^2.
\]
 Hence, as Hermitian forms in
$k,\bar\ell$,
$-R^{p\bar q}{}_{k\bar\ell}(g)\widetilde g_{p\bar q}
\leq C\Lambda g_{k\bar\ell}$.
Contracting with the positive tensor $F^{k\bar\ell}/\Lambda$
bounds the curvature term by $C\mathcal S$.  We have proved
\begin{equation}\label{eq:raw-trace-bound}
 -\mathcal L\log\Lambda
 \leq C(1+\mathcal S)
 +\frac{C}{\Lambda}\tr_{\widetilde g}g
       \bigl(1+\tr_\theta g\bigr).
\end{equation}
Every constant used so far is uniform as $t\downarrow0$.  Notice
that the only inverse of the degenerating form $\theta$ occurs in
the explicitly retained factor $\tr_\theta g$.

It remains to use the divisor weight.  The equation says that the
nonnegative eigenvalues of $\widetilde g^{-1}\theta$ have sum $d$,
so each is at most $d$, i.e. $\theta\leq d\widetilde g$.
Taking inverses reverses the order, and therefore
\[
 \widetilde g^{-1}\leq d\theta^{-1},
 \qquad
 \tr_{\widetilde g}g\leq d\tr_\theta g.
\]
On the other hand, \eqref{eq:B-weighted-hatC} says
$\theta\geq c_1|s|^{2\gamma}g$, whence
$\theta^{-1}\leq c_1^{-1}|s|^{-2\gamma}g^{-1}$.  Taking the trace
against $g$, and using the uniform upper bound for $d=d_t$, gives
\[
 \tr_\theta g\leq C|s|^{-2\gamma},
 \qquad
 \tr_{\widetilde g}g\leq d\tr_\theta g
 \leq C|s|^{-2\gamma}.
\]
Recall that $|s|\leq1$.  Hence the last term in
\eqref{eq:raw-trace-bound} satisfies
\[
 \frac{C}{\Lambda}\tr_{\widetilde g}g
       (1+\tr_\theta g)
 \leq\frac{C}{\Lambda}|s|^{-4\gamma}.
\]
If the right-hand side is larger than $1$, then, after enlarging $C$,
\eqref{eq:weighted-alternative} holds.  Otherwise
$-\mathcal L\log\Lambda\leq C(1+\mathcal S)$.

Finally \eqref{eq:Theta-hatC-w} and the equation give the exact identity
\begin{equation}\label{eq:Lw}
 \mathcal L w_t
 =\mathcal F_t^{i\bar j}(\Theta_t-\widehat C_t)_{i\bar j}
 =d_t-\mathcal S_t.
\end{equation}
At an interior maximum, $\mathcal L Q_t\leq0$. Combining the last two
displays proves \eqref{eq:maximum-inequality}.
\end{proof}

\end{document}
